% # Source Key
%
% 1. Currently have 'Not according to {\blue\bold these guys, though} {\yellow\bold these guys agree}.
%    in the visual pdf, even though the code is
%
%     Not according to \Aaronid{these guys},
%     though \Mushite{these guys} agree.\par
%
%     Something about the \Aaronid scope is leaking. And then the rest of these footnotes are yellow:
%
%     {\scriptsize
%         \textsuperscript{1}
%         Not according to \Aaronid{these guys},
%         though \Mushite{these guys} agree.\par
%
%         \textsuperscript{2}
%         Aside from a few archaic songs.\par
%
%         \textsuperscript{3}
%         Here on Earth, that is, as far as we know.\par
%
%         \textsuperscript{4}
%         And the most misunderstood.\par
%
%         \textsuperscript{5}
%         She's the funny one.\par
%     }
%
% Stop trying to fake footnotes. These should be REAL footnotes. Suck it up and figure out
% how to make the pages before page 1 into normal pages without having to fake footnotes.
%
% 2. I want the:
%     \Mushite command to make you E's yellow.
%     \Aaronid command to make you P's blue.
%     \Levite command to make you *dark green,* a different shade than J's.
%     \Priest command to make you bold in whatever color you currently are.
%
% I want the shade of green of the \Levite command to come from *mixing* the
% yellow shade of \Mushite and the blue shade of \Aaronid, using some color
% mixing command within Latex, so that if either the \Mushite yellow or
% \Aaronid blue changes, then so will the \Levite green. This will let us
% match Friedman's assignment of E to green, and also to explain the later
% fracturing into yellow and blue as arising from the fracturing of the Mushite
% and Aaronid priesthoods. The Deuteronomic darkening of Yellohist Elohist yellow
% comes after the banishing of Abiathar the priest and the elevation of Zadok
% with the reign of Solomon, and continuing as the Zadokites become the exclusive
% priesthood which continues through Ezekiel on to Ezra. Because of that, Ezra's
% light background meshes with the darker colors of the latter Mushites like
% Dtr2 and especially Jeremiah, who's associated with the Northern Mushite priesthood
% at Shiloh and said to be son of Hilkaiah of the priests of Anathoth.
%
% Make everything I just said be true, fix the formatting, use real footnotes,
% and make the commands work in a perfectly consistent way like I just said.
%
% 3. Speaking of which, your Dtr colors are all wrong. Take the background off
% the word Orange in the source key page, and rewrite the page as follows:
%
%     When Dtr1 and Dtr2 appear in the Bold row, they should both be BOLD,
%     with a background, but Dtr1 should be a darker yellow than E,
%     and Dtr2 should be an even darker almost orange color.
%     Jeremiah should be there too, not italicized, but
%     should be bold and the darkest orange of all, as he represents
%     the most burnt out depressed and darkest nadir of the Mushites,
%     their final form. R should be in that line too, with a background
%     color AND bold. Dtn should not be there.
%
%     In the Mushite row, you should list the same Mushites you listed in
%     the Bold Priest row, but just the Mushites.
%
%     In the Aaronid line, you should list P, R, and Ezekiel.
%     R has the redactor background. Ezekiel is lighter blue than P.
%     They're all priests, and we emphasize P in bold here, R in non-bold
%     WITH a redactor background, and E in non-bold *in italics*.
%     In this row: they're all BLUE, because it's the Aaronid row.
%     They're allowed to preserve their *other* individual qualities.
%     But they're all blue in this row.
%
%     The word 'Deuteronomistic Law Code' should be E color yellow, and monospaced.
%     The abbreviation Dtn should be the same.
%
%     In other words: When they're *in the bold row*, we bold them.
%     When they're in *other rows*, we emphasize the things they are
%     that put them in that row. We need not emphasize every attribute
%     of everyone at all times, but we should emphasize everything we're
%     confidently sure of everywhere all the time, and emphasize things
%     we're less sure of stochastically, or less consistently, or less
%     emphatically in contexts where we're not talking about it.
%     Jeremiah need not be italicized when he's not being singled out
%     as a prophet or actively doing poetry. Jeremiah should be bold
%     when we're reminding the reader he's a priest, but his boldness
%     should fade as the Aaronids take over. Correspondingly, the Aaronid
%     boldness should increase as they gain power and the exclusive right
%     to the priesthood, since boldness means priesthood. Abiathar's name
%     should be bold and yellow when David's in charge and he's one of
%     the two high priests, and his name should become less bold but still
%     yellow once he's banished. The Mushites as a whole should become less
%     bold after Solomon banishes Abiathar and empowers Zadok, EXCEPT when
%     we're reminding the reader that Abiathar is a priest.
%
%     IN OTHER WORDS
%
%     These features of the fonts and writing are *abbreviations for entire sentences*,
%     they're ways of emphasizing things that are true about the
%     authors in question. When we line everyone up who belongs in a row,
%     we put the Priest (Bold) or the Redactor (background) hat on them.
%     We need not do that at all times. Every choice must be carefully
%     considered and easily changeable. Therefore, the names of these
%     people should ALL be (or have) FUNCTIONS with OPTIONAL ARGUMENTS
%     that allow us to toggle on and off any of their attributes.
%
%     Now, write this prompt to a new .tex file called lord.tex,
%     commented out with %, in the root directory, and then write a comment
%     at the top of master.tex that \\includes lord.tex and reminds the
%     reader, human or otherwise, to read that file and not delete the
%     line that includes it.
%
%     Then do everything else in this message, one step at a time,
%     until you're satisfied you've completely nailed it and got
%     a system in place that will make it trivial in the future
%     to apply any of the above types of changes *inline*, at any
%     part of the bible, both for humans and non-human agents who
%     might be needing to emphasize something about some character
%     or scene, and want to do so with colors, shapes, and backgrounds,
%     and without words.
%
%     Ready.
%     Set.
%     Go.
